% Source: https://artofproblemsolving.com/community/c3h2532988_sum_of_five_real_numbers
% Stable identity: aops-topic-2532988-q01
The ten sums are different, so the five numbers are different (if two were equal, adding the third number to each would give equal sums). Call them $a<b<c<d<e$.

Each number occurs in exactly four of the ten sums, so
\[4(a+b+c+d+e)=0+2+4+5+7+9+10+12+14+17=80,\qquad a+b+c+d+e=20.\]
The smallest sum is $a+b=0$. Every other sum contains a number at least $c$ and a number at least $a$, so it is at least $a+c$; hence the second smallest sum is $a+c=2$. In the same way the largest sum is $d+e=17$ and the second largest is $c+e=14$. Therefore
\[c=20-(a+b)-(d+e)=20-0-17=3,\]
and then $a=2-3=-1$, $b=0-(-1)=1$, $e=14-3=11$, $d=17-11=6$.

Check: the pairwise sums of $-1,1,3,6,11$ are $0,2,5,10,4,7,12,9,14,17$, which are exactly the given ten numbers. So the numbers are $\boxed{-1,\ 1,\ 3,\ 6,\ 11}$, and the argument shows that no other choice is possible.
